\section{Input for Plasma Background} \label{sec2.5}

\textbf{General remarks}

\medskip

The background medium (mostly plasma) consists of NPLS (sometimes in the code:
NPLSI) different so called ``bulk particle" species, also referred to as ``bulk
ions", by abuse of language.
For each of these species arrays of the parameters: particle density,
temperature, flow velocity (3 Cartesian components),  are defined, one value
per parameter and per grid cell.
There are, in total, NSBOX cells in a run.

From the assumption of quasi-neutrality, the specified charge state of each
background ion species (including  charge state = 0 for neutral background
species), and the data for density $n_i$ of background species labeled $i$,
DIIN(IPLS,icell), an array of NSBOX data for the \textbf{electron density} $n_e
  (\si{\per\cubic\centi\metre} )$ (DEIN(icell) is computed in each of the cells
(see next subsection: ``derived background data").

Further background tallies (or ``input tallies") defined in this block are
related to a magnetic field, electric field, cell volumes (unless computed by
EIRENE defaults).

For each ``background tally" ITALLY there is a flag INDPRO(ITALLY), which
describes the type of profiles, the usage of input parameters, or how a tally
is transferred from an external file or code into an EIRENE run.
The meaning of variable INDPRO is described below.

I.e.: the background medium is described by several blocks of NSBOX data each,
(i.e.\ one datum per grid cell), so called ``input tallies", namely
\begin{itemize}
  \item one array for the electron temperature $T_e$
        (\si{\electronvolt}) (TEIN),
  \item either one common or NPLS distinct (one for
        each bulk ion species) arrays for the ion temperature(s) $T_i$
        (\si{\electronvolt}) (TIIN),
  \item NPLS arrays (one for each bulk ion species),
        ion densities $n_i$ (\si{\per\cubic\centi\metre}) (DIIN),
  \item either one
        common or NPLS distinct (one for each bulk ion species) arrays for the
        (Cartesian) drift velocities $\vV = ( V_x, V_y, V_z )$ (VXIN, VYIN, VZIN),
        either in \si{\centi\metre\per\second} or Mach number units.
\end{itemize}

\minisec{Special case: Magnetic field input}

The unit vector  $\vb$ = $( b_1, b_2, b_3 )$ (BXIN, BYIN, BZIN), parallel to
the magnetic field is set to (0.,0.,1.) by default, i.e.\ pointing in z- (or
toroidal-) direction.
The default magnetic field strength is $|\vB|=1$ T in all grid cells.
This default setting may be overruled by using the \emph{INDPRO(5)} flag and
the corresponding input parameters (see below).
Mostly the full magnetic vector field is read from external files or routines
(\emph{INDPRO(5)} $> 3$ options).
For testing purposes also quite simple preprogrammed magnetic fields can be set
(\emph{INDPRO(5)} $\leq 3$ options).
In this latter case either a (radial or x-) profile of a vector field $\vb$ =
$( 0.0, b_2, b_3 )$ (which is parallel to the $x=$ const.
``flux" surfaces) or, alternatively (since 2016) of a vector field $\vb$ = $(
  b_1, 0.0, b_3 )$ can be defined.
In these two cases a B-field ``pitch"  profile can be selected.

For example, in the first case, ($INDPRO(5) = 1, \dots, 3$, and geometry level
$LEVGEO=1,2,3$), the profile input parameters $B1, \dots, B5$  control the x-
or radial profile $\mathrm{pitch}(x) = b_2/B_{tot}$.
From this a cell centered B-field unit vector is internally constructed, by
assuming that the first set of coordinate surfaces (x- or radial surfaces) are
flux-surfaces (B-field aligned and constant pitch).
The magnetic field direction then has components ($0.
$, pitch, $\sqrt{1-\mathrm{pitch}^2}$) in the three coordinate directions.
1D simulations (x-grid only) are then across the magnetic field.

An analogous procedure is applied in the second case, but there the pitch is
defined as $\mathrm{pitch}(x) = b_1/B_{tot}$ via the input profile parameters
$B1, \dots, B5$, allowing 1D simulations (x grid only) parallel to the magnetic
field (PITCH = 1.0 = const).

Currently only $INDPRO(5)=3$ profile option (piecewise constant profile) allows
to define also the B-field strength (BFIN) [T] together with the field pitch,
solely by the input flags $B1, \dots, B5$ described below.

Alternatively and in all other options an external data source for the B-field
vector in Cartesian coordinates ($INDPRO(5) = 4, \dots, 7$, all geometry levels
$LEVGEO$) is used directly.
In this latter case the transfer of magnetic field data (BXIN,BYIN,BZIN and
BFIN) for each cell into EIRENE proceeds from user supplied profile subroutine
PROUSR ($INDPRO(5)=5$) or, ($INDPRO(5)=4$) from external file or,
($INDPRO(5)=6,7$), via work array RWK, and then from subroutine PROFR).

\medskip

\minisec{Other input tallies: ADIN, ZIIN, VOL, WGHT}
The same applies to the ``additional plasma profile array" ADIN, $INDPRO(6)$,
to permit usage of EIRENE output and graphics routines for additional (derived)
plasma profiles such as pressure, energy fluxes, collision rate coefficients
(see input block 14 for the latter), etc., which are not directly needed in an
EIRENE run but which may be available from an external plasma code, or are
derived from the atomic or molecular data sets included in a particular run.
For all those: $INDPRO(6) = 5, \dots, 7$ in input block 5 only, or the options
in input block 14 (routine AMDIAG.f).

It also applies to the space and species dependent weight function WGHT (to be
used for non-analog sampling techniques, INDPRO(7)), which is defaulted to
1.0D0 for all cells and all NSPZMC test particle species:
NSPZMC = NATM + NMOL + NION + NPHOT.
WGHT is currently unused, and could in principle be eliminated from all current
runs, in particular when storage limitations are a concern.

It furthermore also applies to the zone volume array VOL
\si{\per\cubic\centi\metre}, $INDPRO(12)$.
Defaults are the cell volumes computed in EIRENE subroutine VOLUME from the
standard mesh data.

Additionally it also applies to the ion charge array ZIIN, $INDPRO(11)$.
Defaults are the charge states of the respective bulk ions.

Plasma data in the additional zones IADD, IADD=NSURF+1,NSURF+NADD outside the
standard mesh are defaulted to the ``EIRENE vacuum data" given below.
Plasma data other than these ``vacuum values" (meaning: no collision processes
in such cells) in these zones \textsl{have} to be specified either explicitly
in input block 8 (see below), or in the user supplied subroutine PLAUSR or with
the INDPRO=7 option (see below, this section).

Cell volumes in the additional cell region are defaulted to
\SI{1.0}{\per\cubic\centi\metre}, unless $INDPRO(12)=4, 5, 6, 7$.

\minisec{Vacuum zones: no collision processes}
A standard mesh zone ICELL is automatically identified as a vacuum zone with
regard to plasma species IPL (no collisions with the bulk particles IPL there,
\emph{LGVAC(ICELL,IPL)=.TRUE.}, for electrons: IPL=NPLS+1)

if at least one of the following conditions is met:

for electrons: IPL=NPLS+1:

TEIN(ICELL) $\le$ TVAC,

DEIN(ICELL) $\le$ DVAC

for bulk ion species IPL=IPLS:

TIIN(ICELL,IPL) $\le$ TVAC,

DIIN(ICELL,IPL) $\le$ DVAC

and for all background species, IPL=0:

LGVAC(ICELL,I)=TRUE for I=1,NPLSI and for I=NPLS+1.

The EIRENE vacuum data are defaulted (in subroutine PLASMA) to:

DVAC = 1.0000E 02 \si{\per\cubic\centi\metre}

TVAC = 2.0000E-02 \si{\electronvolt}

Furthermore, zero plasma drift velocities are set in cells which are treated as
vacuum zones.

VVAC = 0.0000E 00  \si{\centi\metre\per\second}

for all three Cartesian components and for all NPLS background particle
species.

Note: A cell can be considered a vacuum cell with respect to electrons,
(LGVAC(\dots,NPLS+1) = TRUE), without being a vacuum cell for all background
ions.
This is because, by abuse of language, also neutral particle species may be
used as background ``ion" species (NCHARP = 0), to include neutral--neutral
collisions, e.g., by the iterative option NITER (input block 1).

\medskip

\textbf{The Input Block}

\begin{lstlisting}
*** 5. Data for Plasma background
      NPLSI
      DO 51  J= 1, NPLSI
\end{lstlisting}

\textsl{* read NPLSI species blocks with \$ = P}

\begin{lstlisting}
        ISPZ$ TEXT$ NMASS$ NCHAR$ NPRT$ NCHRG$ ISRF$ ISRT$
     .    ID1$
        NRC$  NFOL$ NGEN$  NHSTS$ ID3$
C      (format: I2,1X,A8,1X,12(I2,1X))

C End of NPLSI species cards reading
   51 CONTINUE

      (INDPRO(J), J=1,12)
      IF (INDPRO(1).LE.5) THEN
        TE0  TE1  TE2  TE3  TE4  TE5
      ENDIF
      IF (INDPRO(2).LE.5) THEN
        (TI0(I) TI1(I) TI2(I) TI3(I) TI4(I) TI5(I) I=1,NPLSI)
      ENDIF
      IF (INDPRO(3).LE.5) THEN
        (DI0(I) DI1(I) DI2(I) DI3(I) DI4(I) DI5(I) I=1,NPLSI)
      ENDIF
      IF (INDPRO(4).LE.5) THEN
        (VX0(I) VX1(I) VX2(I) VX3(I) VX4(I) VX5(I) I=1,NPLSI)
        (VY0(I) VY1(I) VY2(I) VY3(I) VY4(I) VY5(I) I=1,NPLSI)
        (VZ0(I) VZ1(I) VZ2(I) VZ3(I) VZ4(I) VZ5(I) I=1,NPLSI)
      ENDIF
      IF (INDPRO(5).LE.5) THEN
        B0  B1  B2  B3  B4  B5
      ENDIF
      IF (INDPRO(11).NE.0.AND.INDPRO(11).LE.5) THEN
        (ZI0(I),ZI1(I),ZI2(I),ZI3(I),ZI4(I),ZI5(I) I=1,NPLSI)
      END IF
      IF (INDPRO(12).LE.5) THEN
        VL0  VL1  VL2  VL3  VL4  VL5
      ENDIF
C     ``Option-Card'' for additional information for input tallies
C     only for EIRENE versions starting from
C     Eirene_revised_input_tallies
\end{lstlisting}

If an ``Option-Card'' is encountered then
read an arbitrary number ``input tally option cards'' of the format

\begin{lstlisting}
      ITAL IOPT
\end{lstlisting}

Currently option INDPRO=6 and INDPRO=7 are not available for cell volumes
(input tally no.\ 12).

\medskip

\textbf{Meaning of the Input Variables for Plasma Parameters}

\medskip

The meaning of the variables in the bulk ion species cards is as in block 4
(\cref{sec2.4}) for the test particle species cards specified there.

\textbf{New options since 2003:}

However, after the ID3 flag, there may be 2 additional input flags: (only for
Eirene$_{2003}$ and younger):

\begin{lstlisting}
... CDENMODEL NRE
\end{lstlisting}

Format: \dots,1X,A10,1X,I2.

These flags, if included, control additional options to set the density,
temperature and flow field of the selected species IPLS.
If the string CDENMODEL is found here for a species IPLS, then, after reading
the species and reaction cards for IPLS, one (default) or NRE further input
cards are expected.
\begin{description}
  \item[CDENMODEL] ~

        \begin{description}
          \item[= 'Multiply'] docu to be written
          \item[= 'Constant']
                docu to be written
          \item[= 'fort.13'] Read data for this species from file
                fort.13.

                The additional card (only one) reads (format 2I6):

                \begin{lstlisting}
      ISP ITP
\end{lstlisting}

                ISP is the number of a bulk species on the file fort.13 (e.g. written in an
                earlier EIRENE run, see NFILEL option in input block 1).
                The parameters (fields of density, temperature, flow velocity) for species IPLS
                are set from those of species ISP on fort.13

                ITP is the type of the particle, i.e., always: ITP=4 for this option.
                Internally set: ITP=4, bulk particle type, independent of input value for ITP.

          \item[= 'fort.10'] Read data for this species from file fort.10.

                The additional card (only one) reads (format 3I6):

                \begin{lstlisting}
      ISP ITP ISTR
\end{lstlisting}

                ISP, ITP is the number and type of a test particle species on the file fort.10,
                respectively, from stratum no.\ ISTR (i.e.\ written onto fort.10 either in the
                present run or an earlier EIRENE run, see NFILEN option in input block 1).
                The parameters (fields of density, temperature, flow velocity) for species IPLS
                are derived from those of species ISP, ITP, ISTR on fort.10.

                Note: at the end of a full EIRENE run the corresponding bulk particle profiles
                are reset to those evaluated in this run, for further use in post processing
                routines such as the diagnostic module (see \cref{sec2.11}), printout, plotting
                or iterative mode (input flags NITERI, NITERE, and user subroutine MODUSR, see
                \cref{sec2.1})

          \item[= 'Saha'] Set parameters from Saha-equilibrium of ionization states.
                The additional card (only one) reads:

                \begin{lstlisting}
      ISP ITP ISTR ...
\end{lstlisting}

                to be written
          \item[= 'Boltzmann']
                Set parameters
                from Boltzmann-equilibrium of excited states.
                The additional card (only one) reads (format: 3I6,6x,2E12.4)

                \begin{lstlisting}
      ISP ITP ISTR G_BOLTZMANN DELTA_E
\end{lstlisting}

                docu to be written
          \item[= 'Planck'] (only for background photons) to be
                written
          \item[= 'Corona'] Set density parameters from Corona-equilibrium to a
                ``donor species", using an excitation rate coefficient from atomic database
                FILNAM and a fixed radiative decay rate $A=A_{CORONA}$.
                The additional card (only one) reads (format: 3I6, 1X, A6, 1X, A4, A9, A3,
                E12.4):

                \begin{lstlisting}
      ISP ITP ISTR FILNAM H123 REACTION CR A_CORONA
\end{lstlisting}

                ISP, ITP is the number and type of a (donor) particle species, either one of
                the already defined bulk species IPLS: ISP=IPLS, (for ITP=4), or from the file
                fort.10 (for ITP=0,1,2 or 3) from stratum no.\ ISTR (i.e. written onto fort.10
                either in the present run or an earlier EIRENE run, see NFILEN option in input
                block 1).
                I.e.\ the density $ni(IPLS)$ at each position $x$ (grid cell) for this
                background species IPLS is specified via:
                \begin{equation}
                  \nonumber
                  ni(IPLS,x)= n_{IS}
                  (x) \times \frac{Rate_{IS}(x)}{A}
                \end{equation}
                Here $Rate_{IS}$ is the rate
                for excitation from species $IS$ to $IPLS$ from the A\&M data file $FILNAM,
                  \dots$, and $A$ is a radiative decay rate $A = A_{CORONA}$ read from this input
                card.
                The parameters (fields of density, temperature, flow velocity) for donor
                species IS are taken to be those of species [ISP, ITP, ISTR].

                I.e.: this model constructs a background species IPLS which is in
                Corona-equilibrium with these IS = [ISP,ITP, ISTR] donor particles.
          \item[= 'Colrad'] set parameters density, temperature and flow velocity from
                collisional radiative equilibrium with one or more (NRE) ``donor-species",
                using reduced population coefficients from atomic database FILNAM.
                I.e.\ the density ni(IPLS) at each position $x$ (grid cell) for this background
                species IPLS is specified via:
                \begin{equation}
                  \nonumber
                  ni(IPLS,x)= \sum_{IS} n_{IS}
                  (x) \times POP_{IS}(x)
                \end{equation}
                The NRE (for NRE different donor-states)
                additional cards read:

                \begin{lstlisting}
      ISP ITP ISTR FILNAM H123 REACTION CR
\end{lstlisting}

                ISP, ITP is the number and type of a particle (donor) species $IS$, either one
                of the already defined bulk species IPLS: ISP=IPLS, (for ITP=4, then ISTR is
                irrelevant), or from the file fort.10 (for ITP=0,1,2 or 3) from stratum no.\
                ISTR (i.e. written onto fort.10 either in the present run or an earlier EIRENE
                run, see NFILEN option in input block 1).
                The parameters (fields of density, temperature, flow velocity) for donor
                species IS are set from those of isotope [ISP, ITP, ISTR].
                And by multiplication with a reduced population coefficient $POP$ specified via
                $[FILNAM, H123, \dots]$ these are turned into excited state or other (ionised
                state) populations.
                The $POP$ coefficients are typically from externally applied collisional
                radiative models.

                I.e.\ this model constructs a background species IPLS which is the sum of NRE
                species, each of which being in local collision-radiative-equilibrium with its
                donor isotope $IS$.
        \end{description}
\end{description}

\begin{description}
  \item[INDPRO] Flag-array for the type of profile.
        The last digit (between 1 and 9) controls the type of profile.
        A second digit and/or the sign control further options, as described below.
        A third digit for each entry can be used to switch from a cell-wise constant
        profile (default) to a smooth (interpolated into cells) profile.
        This option is currently being tested for the magnetic field input profiles.
        I.e.\ INDPRO(5)=103 provides the same magnetic field as INDPRO(3)=3, however
        the magnetic field is evaluated on the fly (along particle trajectories) by
        interpolation, at the very point of particle position rather than at the cell
        center (=const.
        per cell).

        Options INDPRO = 1, 2, 3 are preprogrammed 1 D profiles with the first
        coordinate (x or radial) taken as independent parameter.
        I.e.\ profiles are constant in the other 2 coordinate directions: y and z.

        Options INDPRO = 4, 5, 6, \dots\ allow more general, and higher dimensional
        profile forms.

        INDPRO is an array of length 12.
        Each element in this array controls one particular input tally, namely:
        \begin{description}
          \item[INDPRO(1)]  for TEIN

          \item[INDPRO(2)]  for TIIN By the default ($\emph{0} < \emph{INDPRO(2)} <
                  \emph{10}$): one common $T_i$ profile for all ``bulk ion" species is set.
                I.e., read only one profile card.

                \textbf{New options since 2001:}

                Values of INDPRO(2) larger than 10: use only the last digit, and one $T_i$
                profile card must be read for each bulk ion species IPLS.
                For example: \emph{INDPRO(2)=15} or \emph{=25}, means: one separate ion
                temperature must be specified for each bulk ion species, and the profile type
                is 5.
                (\emph{INDPRO(2)=-5} would do the same.)
          \item[INDPRO(3)] read NPLSI cards, for DIIN, IPLS = 1,NPLSI

          \item[INDPRO(4)] read NPLSI cards, for VXIN, VYIN, VZIN, IPLS = 1,NPLSI

                \textbf{New options since 2001:}

                By the default ($\emph{0} < \emph{INDPRO(4)} < \emph{10}$): one separate flow
                field for each bulk species is set, velocity is given in
                \si{\centi\metre\per\second}.

                Negative value of INDPRO(4) means: Mach number units instead.
                The sound speed \emph{cs} is taken to be the isothermal ion acoustic speed of
                species IPLS.

                \emph{ABS(INDPRO(4)} is then used as flag for the choice of profile type.

                Values of INDPRO(4) larger than 10: use only the last digit, and only one
                common flow field is set for all bulk ion species.
                Note the difference to the $T_i$ options:
                there the meaning of INDPRO larger than 10 was exactly opposite
                to the meaning here
                for the flow fields (due to historical reasons and for backward
                compatibility of input
                files.
                EIRENE had originally by default one single common ion temperature, but one
                flow field for each bulk ion species).

          \item[INDPRO(5)] for a 1D (radial or x-) profile of PITCH (later, in
                initialization phase,  converted into Cartesian unit B-field vector
                BXIN,BYIN,BZIN)

                By the default ($\emph{0} < \emph{INDPRO(5)} < \emph{10}$) the pitch is defined
                as $B_y/B_{tot}$ (LEVGEO=1), $B_{\theta}/B_{tot}$ (LEVGEO=2), or as
                $B_{pol}/B_{tot}$ (LEVGEO=3), where $B_{pol}$ is the direction along the
                polygons.

                Negative values of INDPRO(5) (only for LEVGEO=1): use \emph{ABS(INDPRO(5)}, and
                the pitch is then defined as $B_x/B_{tot}$ (only for LEVGEO=1), allowing 1D
                simulations, using the first (x-) grid only, along the B field.
                This option was only added during 2016.

                \textbf{New options since 2001:} In case INDPRO(5)=3 (flat profile) the two
                redundant input parameters B2, B3 are used to define a constant B-field
                strength [T], see below under profile type INDPRO=3.

          \item[INDPRO(6)] for ADIN

          \item[INDPRO(7)] for WGHT (to be written)

          \item[INDPRO(12)] for VOL.

        \end{description}

        For the input tally profiles no.\ 6 and 7 (ADIN, WGHT) only the options
        INDPRO=5 or INDPRO=6 exist.

        For the profile no.\ 12 (cell volumes) the options INDPRO(12) = 4, 5, 6, or
        INDPRO(12) = 7 are active options.
        For all other values of INDPRO for these latter four profiles the default
        profiles described above (``General remarks") are set.

        For each profile up to 6 parameters P0, \dots, P5 are read, e.g.\ TE0, \dots,
        TE5 for the $T_e $-profile, TI0, \dots, TI5 for the $T_i $-profile, and so on.

        Depending upon the value of INDPRO one of the profile routines PROFN, PROFE,
        PROFS, etc.\ is called from subroutine PLASMA.
        \begin{description}
          \item[INDPRO = 1-4] ~

                NR1STM plasma data are defined on the one- dimensional zone-centered grid

                RHOZNE(J), J=1, \ldots,NR1STM	(NR1STM = NR1ST - 1 ) .

                Vacuum data are specified in zone NR1ST.

                These \hfill NR1ST \hfill data \hfill are \hfill copied \hfill NBLCKS \hfill
                times \hfill to \hfill define

                NBLCKS identical profiles, totally : NBLCKS $\cdot$ NR1ST = NSURF data
                (subroutine MULTI).

                Here the number of copies is NBLCKS = NMULT $\cdot$ NP2ND $\cdot$ NT3RD.
                see ``Input Data for Standard Mesh" (block 2).
          \item[INDPRO = 1] ~

                (see subroutine PROFN)

                Parameter P5 must be a position inside the first (radial) grid

                RHOSRF(1) $\le$ x $\le$ P5:

                $P(x) = P1+( P0 - P1 ) \cdot {\left ( 1 - {\left ( \frac{x - RHOSRF(1)}{P5 - RHOSRF(1)} \right )}^{P2} \right )}^{P3}$

                \medskip

                in particular :
                \[
                  \begin{array}{l}
                    x = RHOSRF(1)   \rightarrow  P(x) = P0 \\
                    x = P5	    \rightarrow  P(x) = P1
                  \end{array}
                \]

                \medskip

                P5 $\le$ x $\le$ RHOSRF(NR1ST) :

                P(x) = P1  $\cdot$ exp(- (x - P5) / P4)
          \item[INDPRO = 2] ~

                (see subroutine PROFE)

                Parameter P5 must be a position inside the first (radial) grid

                RHOSRF(1) $\le$ x $\le$ P5 :

                P(x) = P0  $\cdot \exp$(- (x - P1) / P2)

                in particular : x = P1 $\rightarrow$ P(x) = P0

                P5 $\le$ x $\le$ RHOSRF(NR1ST) :

                P(x) = P(P5)  $\cdot \exp$(- (x - P5) / P4)

                (P3 : no meaning)
          \item[INDPRO = 3] ~

                (see subroutine PROFS, PROFS(P0,P1,P5,PVAC)

                Parameter P5 must be a position inside the first (radial) grid

                RHOSRF(1) $\le$ x $\le$ P5 :	   P(x) = P0

                P5 $\le$ x $\le$ RHOSRF(NR1ST) :    P(x) = P1

                Parameters P2, P3, P4 have no meaning, except for the magnetic-field (INDPRO(5)
                = 3): Here PROFS(P0, P1, P5, PVAC) is used to define the pitch angle profile,
                and it is also used to set the B-field strength (\si{\tesla}) by calling
                PROFS(P2, P3, P5, PVAC)

          \item[INDPRO = 4] ~

                read from input stream fort.P0

          \item[INDPRO = 5] ~

                EIRENE calls a user supplied plasma profile routine:

                Subroutine PROUSR (RHO, INDEX, P0, P1, P2, P3, P4, P5, PVAC, NDAT).

                NDAT data must be defined on

                RHO(J), J = 1, NDAT e.g.\ by using the profile parameters P0, \dots, P5 and any
                user supplied profile function.
                See \cref{sec3.6}.

                (NDAT=NSURF)

                The flag INDEX is as for the ``INDPRO = 6" option, see below, and
                \cref{sec3.6,sec4}.
                It determines, which of the background medium profiles has to be provided at a
                particular call to subroutine PROUSR.

                PVAC is the EIRENE vacuum value for the profile in question.
          \item[INDPRO = 6] ~

                EIRENE calls a the routine:

                Subroutine PROFR (RHO,INDEX,N1I,N1,NDAT).

                N1I (e.g., NPLSI) profiles, of NDAT = NSURF data each, are read onto RHO(I1I,
                J), J = 1, NDAT, I1I = 1, N1I from a work array RWK.
                N1 (e.g., NPLS) is the leading dimension of the array PRO as specified in the
                calling program.
                The data must be written onto array RWK in the initialization phase, e.g., in
                subroutine INFCOP or MODUSR.
                The location of a particular profile on this work array is determined by the
                value of INDEX, see \cref{sec4} below.
                By this option plasma parameters may directly be transferred into EIRENE from
                other files, e.g.\ from plasma transport codes.

                NDAT (= NSURF) is the number of cells in the ``Standard Mesh".
                Parameters in the additional cell region

                ICELL = NSURF + 1, \dots, NSURF + NRADD

                are set to the ``EIRENE vacuum values"

                For further details of subroutine INFCOP see \cref{sec2.1} (``Input Data for
                operating mode"), input variable NMODE and \cref{sec2.14,sec4} for an example.

          \item[INDPRO = 7] ~

                Same as INDPRO = 6 option, except that NDAT = NSBOX rather than NDAT = NSURF.
                Hence in this case background data are read from RWK also for the additional
                cell region

                ICELL=NSURF+1,NSURF+NRADD.

          \item[INDPRO = 8] ~

                to be written

        \end{description}
\end{description}

For EIRENE versions starting from Eirene\_revised\_input\_tallies additional
input for use with input tallies is available.

\begin{description}
  \item[``Option-Card''] If EIRENE finds a card containing the string ``OPTIONAL''
        it assumes that there might be an arbitrary number of input cards
        describing additional options for input tallies.
  \item[ITAL] index of input tally, -NTALI < ITAL < 0
  \item[IOPT] option flag
        \begin{description}
          \item[$< 0$] switch off tally ITAL
          \item[$= 0$] keep default
          \item[$= 1$] explicitely switch on tally ITAL
          \item[$= 2$] prepare for smoothing, interpolate tally ITAL to cornerpoints
          \item[$= 3$] switch on gradients for tally ITAL
        \end{description}
\end{description}

\subsection{Derived Background Data} \label{sec2.5.1}

Given the data describing the background medium (plasma) in each cell (on
common block COMUSR) a set of further, ``derived background data" profiles is
computed and stored also on COMUSR.
This is done by a call to subroutine \lstinline|PLASMA_DERIV|.
The following quantities are defined there:
\begin{description}
  \item[DEIN] Electron density (\si{\per\cubic\centi\metre}), derived from all
        NPLSI ion densities and the charge states NCHRGP(IPLS) of each ion species.
  \item[EDRIFT] Kinetic energy (\si{\electronvolt}) in drift motion, for each
        background species IPLS, derived from the mass MASSP(IPLS) and the flow
        velocity VXIN,VYIN,VZIN.
  \item[LGVAC] Vacuum flag (logical) to identify regions, in which no collision
        data for all or certain background species are computed, i.e., the
        collisionalities are set to zero in these cells for those background species.
        See also explanation above, with respect to the EIRENE vacuum data
        TVAC,DVAC,VVAC.
  \item[NSTGRD] flag for special treatment of some selected cells
        \begin{description}
          \item[$= 0$] default
          \item[$= 1$] This cell is a dead cell,
                not to be seen by the particles (isolated from the computational domain.
                NSTGRD(ICELL)=1 can be specified in problem specific routines (\dots USR) or in
                routines interfacing EIRENE with external codes: INFCOP (code segment:
                COUPLE\dots)
          \item[$= 2$] indicates that this cell belongs to a grid cut, i.e.,
                is not a valid cell for particle tracing.
          \item[$= 3$] indicates that this cell is not a real cell but only the storage
                position for spatially averaged tallies.
                E.g.: cell no.\ NR1ST, 2*NR1ST, etc.
        \end{description}
\end{description}
